\documentclass{report}
\usepackage{amsmath,amssymb}
\setlength{\parindent}{0mm}
\setlength{\parskip}{1em}
\begin{document}
\begin{center}
\rule{6in}{1pt} \
{\large Johnathan M. Bardsley \\
{\bf Sampling from Bayesian Inverse Problems with L1-type Priors using Randomize-then-Optimize}}

Math Sciences \\ University of Montana \\ 32 Campus Drive
\\
{\tt bardsleyj@mso.umt.edu}\\
Zheng Wang\\
Cui Tiangang\\
Youssef Marzouk\\
Antti Solonen\end{center}

We focus on applications arising in inverse problems in which the
measurement model has the form
$$
y = F(\theta)+\epsilon,
$$
where $y$ is the measurement vector; $F$ is the forward model function
with unknown parameters $\theta$; and $\epsilon$ denotes independent and
identically distributed
Gaussian measurement error, i.e., $\epsilon\sim N(0,\sigma^2 I)$.
Then the probability density function for the measurements $y$ given the
unknown parameters $\theta$ is given by
$$
p(y|\theta)\propto\exp\left(-\frac{1}{2\sigma^2}\Vert y-F(\theta)\Vert^2_2\right),
$$
where `$\propto$' denotes proportionality. In Bayesian inverse problems,
one also assumes a {\em prior} probability density function $p(\theta)$,
which incorporates both prior knowledge and uncertainty about the unknown
parameters $\theta$. In this talk, we focus on the case in which
the prior $p(\theta)$ is of L1-type, i.e.,
$$
p(\theta)\propto \exp\left(-\lambda \Vert D\theta\Vert_1\right).
$$
Such priors include the total variation prior and the Besov $B_{1,1}^s\;$ space priors.
With these two probability models ($p(y|\theta)$ and $p(\theta)$) in
hand, by Bayes' Law, the {\em posterior} density function has the form
\begin{eqnarray*}
p(\theta|y)&\propto& p(y|\theta)p(\theta)\\
&\propto& \exp\left(-\frac{1}{2\sigma^2}\Vert
y-F(\theta)\Vert^2_2-\lambda \Vert D\theta\Vert_1\right).
\end{eqnarray*}
Regardless of the form of $F$, $p(\theta|y)$ is non-Gaussian due to the
presence of the L1-norm. Moreover, in inverse problems, $\theta$ is
high-dimensional. Taken together,
these challenges make the problem of sampling from $p(\theta|y)$ -- which
is a requirement if one wants to perform uncertainty quantification --
difficult.
To overcome this, we extend the Randomize-then-Optimize (RTO) method,
which was recently developed for posterior sampling when $F$ above is
nonlinear and $p(\theta)$ is Gaussian.
The extension of RTO to the L1-type prior case requires a variable
transformation, which turns $p(\theta)$ into a Gaussian probability
density
in the transformed variables, thus allowing for the application of RTO.
In this talk, we will begin by presenting the RTO method, and then its
extension to the L1-type prior case via the variable transformation.
Several numerical experiments will also be presented to illustrate the
approach and the resulting Markov Chain Monte Carlo method.


\end{document}
